\textbf{Task restriction improves proxy accuracy with calibration and exposure tradeoffs.} \textbf{a--c}, All ten prior candidates on the Kumo-small proxy after 10,000 cellular tasks, 625 updates and up to 2k HVGs. Open circles are three paired task-stream replicates from one model seed; diamonds are means. Gray/blue ticks mark the effective noise band/keep threshold. Positive differences favour each candidate against its specified base. Accuracy endpoints differ across candidates; donor AUC and NLL reduction use separate axes. Supplementary Table~\ref{tab:s5} defines bases and endpoints. Routing-only labels are not efficacy decisions; A8d is not testable because proxy headroom is insufficient. \textbf{d,e}, CellPFN continuation over 4,000 updates and 16,000 tasks per arm. Standard training supplied 12,816 metadata, 2,398 cell-type and 786 donor-identity tasks; restriction supplied 12,001 cell-type and 3,999 donor-identity tasks and excluded signed-covariance banks. Cell exposures were 6,386,352 versus 10,671,056 (1.67-fold), including reuse. \textbf{f,g}, Restricted minus standard balanced accuracy and NLL at update 100,000 with all genes: paired means and 95\% $t$ intervals over twenty support draws. Positive NLL change is worse. Task composition, bank acceptance and cell exposure change together; intervals do not cover training-seed variation. pp, percentage points. Source Data retain full frozen decisions and numerical summaries.
